\section{The Hrushovski--Weil theorem}
\label{sec:hw}

\subsection{Generic sets of germs}

\begin{lemma}[independent generic tuples are conjugate]\label{lem:indep-conj}
Let $V_1,\dots,V_r$ be irreducible $k$-varieties and let $(y_1,\dots,y_r)$,
$(y'_1,\dots,y'_r)$ be two independent tuples ($\rk(y_1\cdots y_r)=\sum_i\dim V_i$,
likewise for the primed tuple) with $y_i,y'_i$ generic points of $V_i$ over $k$.  Then
there is $\sigma\in\Gamma$ with $\sigma y_i=y'_i$ for all $i$.  The same holds with
$\Gamma$ replaced by $\Gamma_M$ for a small algebraically closed $M$ when all the tuples
are generic over $M$ and independent over $M$ (same proof with $k$ replaced by $M$).
\end{lemma}
\begin{proof}
Induction on $r$.  For $r=1$ this is \cref{lem:generic-char}(2).  Given $\sigma_1\in\Gamma$
with $\sigma_1(y_1,\dots,y_{r-1})=(y'_1,\dots,y'_{r-1})$, the tuples $\sigma_1y_r$ and
$y'_r$ are both generic points of $V_r$ over $\acl(y'_1\cdots y'_{r-1})$ (independence
and \cref{lem:generic-char}(1)), hence conjugate under
$\Gamma_{\acl(y'_1\cdots y'_{r-1})}$ by $\tau$; take $\sigma=\tau\sigma_1$.
\end{proof}

Throughout this section $V$ is an irreducible $k$-variety of dimension $m$ and
$V_{\Gen}$ is the set of generic points of $V$ over $k$ (\cref{def:dom}); germs in
tuple positions stand for their codes (\cref{conv:germ-tuple}).

\begin{definition}[generic set of germs]\label{def:genset}
A \emph{generic set of germs} on $V$ is a nonempty subset
$\mathcal G\subseteq\Germ^\times(V)$ which is a single $\Gamma$-orbit.  All its members
have the same rank, denoted $n=\rk(\mathcal G)$, and their code tuples are the generic
points of one irreducible $k$-variety $W=\Loc(\code g)$ (\cref{rem:code-param}), and
$g\mapsto\code g$ is a bijection from $\mathcal G$ onto $W_{\Gen}$.
\end{definition}

\begin{lemma}\label{lem:genset-basic}
Let $\mathcal G$ be a generic set of germs on $V$ of rank $n$.
\begin{enumerate}[nosep]
\item For every finite tuple $w$ there is $h\in\mathcal G$ with $h\ind{}w$; any two such
  $h,h'$ are $\Gamma_{\acl(w)}$-conjugate.
\item If $g\in\mathcal G$ and $y\in V_{\Gen}$ with $y\ind{}g$, then $y\in\operatorname{Dom}(g)$
  and $g(y)\in V_{\Gen}$ with $g(y)\ind{}g$.
\item $\mathcal G^{-1}=\{g^{-1}:g\in\mathcal G\}$ is a generic set of germs of rank $n$.
\end{enumerate}
\end{lemma}
\begin{proof}
(1) The codes of elements of $\mathcal G$ are the generic points of $W$ over $k$ and
determine the germ (\cref{rem:code-param}); take a generic point of $W$ over $\acl(w)$
(\cref{lem:genexists}) and the germ with that code.  Two such codes are generic over
$\acl(w)$, hence conjugate under $\Gamma_{\acl(w)}$ (\cref{lem:generic-char}(2)), and
so are the germs (\cref{lem:germ-aut}).
(2) \cref{lem:eval}(2) with $F=\acl(\Cb g)$.  (3) $\Cb(g^{-1})=\Cb(g)$ and
$\sigma(g^{-1})=(\sigma g)^{-1}$ (\cref{lem:inverse}).
\end{proof}

\subsection{Statement}

\begin{definition}[Hrushovski--Weil hypotheses]\label{def:hw-hyp}
A generic set of germs $\mathcal G$ on $V$ is \emph{generically transitive} if
\begin{enumerate}[label=\textup{(H\arabic*)}]
\item\label{H:trans} for all $g\in\mathcal G$ and $y\in V_{\Gen}$ with $y\ind{}g$:
  $g(y)\ind{}y$;
\end{enumerate}
it is \emph{closed under inverse and generic composition} if
\begin{enumerate}[label=\textup{(H\arabic*)},start=2]
\item\label{H:inv} $g\in\mathcal G\Rightarrow g^{-1}\in\mathcal G$;
\item\label{H:comp} for all $g_1,g_2\in\mathcal G$ with $g_1\ind{}g_2$: $g_1g_2\in\mathcal G$.
\end{enumerate}
\end{definition}

\begin{lemma}[products of independent generics are generic over each factor]
\label{lem:prod-indep}
Let $\mathcal G$ be a generic set of germs of rank $n$, and $g_1,g_2\in\mathcal G$ with
$g_1\ind{}g_2$ and $g_1g_2\in\mathcal G$.  Then $g_1g_2\ind{}g_1$ and $g_1g_2\ind{}g_2$.
\end{lemma}
\begin{proof}
By \cref{lem:comp-codes}(1), $g_1g_2$ and $g_2$ are interalgebraic over $g_1$, so
$\rk(g_1g_2/g_1)=\rk(g_2/g_1)=\rk(g_2)=n=\rk(g_1g_2)$; symmetrically
$\rk(g_1g_2/g_2)=\rk(g_1/g_2)=n$.
\end{proof}

\begin{theorem}[Hrushovski--Weil for germs]\label{thm:hw}
Let $\mathcal G$ be a generically transitive generic set of germs on $V$ of rank $n$,
closed under inverse and generic composition, and let $G\le\Germ^\times(V)$ be the
subgroup generated by $\mathcal G$.  Then:
\begin{enumerate}[label=\textup{(\alph*)}]
\item\label{hw:two} $G=\mathcal G\cdot\mathcal G=\{g_1g_2:g_1,g_2\in\mathcal G\}$, and $G$ is
  $\Gamma$-invariant.
\item\label{hw:stab} (generic stabilisation) for $\gamma\in G$ and $h\in\mathcal G$ with
  $h\ind{}\gamma$: $\gamma h\in\mathcal G$, $h\gamma\in\mathcal G$, $\gamma h\ind{}\gamma$ and
  $h\gamma\ind{}\gamma$.
\item\label{hw:rank} every $\gamma\in G$ has $\rk(\gamma)\le n$, with equality iff
  $\gamma\in\mathcal G$.  In particular $\mathcal G$ is the set of elements of $G$ of maximal
  rank and it is a single $\Gamma$-orbit (``$G$ is connected with generic type
  $\mathcal G$'').
\item\label{hw:action} (generic action) for $\gamma\in G$ and $y\in V_{\Gen}$ with
  $y\ind{}\gamma$: $\gamma(y)\in V_{\Gen}$ and $\gamma(y)\ind{}\gamma$; and
  $(\gamma\delta)(y)=\gamma(\delta(y))$ whenever $y\ind{}\gamma\delta$.
\item\label{hw:space} (homogeneous space) there is a set $S$ with a transitive and
  faithful action of $G$, and an injection $\iota:V_{\Gen}\to S$ such that
  $\gamma\cdot\iota(y)=\iota(\gamma(y))$ whenever $y\ind{}\gamma$; every element of $S$ is
  of the form $\gamma\cdot\iota(y)$.  Explicitly $S=(G\times V_{\Gen})/E$ where
  \[
    (\gamma,y)\,E\,(\gamma',y')\iff (h\gamma)(y)=(h\gamma')(y')
    \text{ for some (equivalently every) } h\in\mathcal G\text{ with } h\ind{}\gamma\gamma'yy',
  \]
  $\delta\cdot[\gamma,y]=[\delta\gamma,y]$ and $\iota(y)=[\mathrm{id},y]$.
\end{enumerate}
\end{theorem}

\subsection{Proof}

We use freely that for $\delta,\gamma\in\Germ^\times(V)$: $\Cb(\delta\gamma)\subseteq
\acl(\delta,\gamma)$ and $\Cb(\delta)\subseteq\acl(\delta\gamma,\gamma)$ (\cref{lem:comp-codes}(1)),
so that \emph{$\delta\gamma$ and $\delta$ are interalgebraic over $\gamma$}; consequently
$\rk(\delta\gamma/C\gamma)=\rk(\delta/C\gamma)$ for every $C$.  Whenever \ref{H:comp} is
applied, \cref{lem:prod-indep} gives in addition that the product is independent from
each factor.

\begin{lemma}[three generators reduce to two]\label{lem:threetwo}
For $g_1,g_2,g_3\in\mathcal G$ there are $g_5,g_6\in\mathcal G$ with $g_1g_2g_3=g_5g_6$.
\end{lemma}
\begin{proof}
Choose $g_0\in\mathcal G$ with $g_0\ind{}g_1g_2g_3$ (\cref{lem:genset-basic}(1)).  Then
$g_1g_2g_3=(g_1g_0)(g_0^{-1}g_2g_3)$.  As $g_0^{-1}\in\mathcal G$ (\ref{H:inv}) and
$g_0^{-1}\ind{}g_2$, \ref{H:comp} and \cref{lem:prod-indep} give $g'':=g_0^{-1}g_2\in\mathcal G$
with $g''\ind{}g_2$.
Since $g''$ and $g_0$ are interalgebraic over $g_2$,
$\rk(g''/g_2g_3)=\rk(g_0/g_2g_3)=n$, so $g''\ind{}g_3$ and $g''g_3\in\mathcal G$ by
\ref{H:comp}.  Also $g_1\ind{}g_0$ gives $g_1g_0\in\mathcal G$.  Take $g_5=g_1g_0$,
$g_6=g''g_3$.
\end{proof}

\begin{proof}[Proof of \cref{thm:hw}\ref{hw:two}]
Let $P=\mathcal G\mathcal G$.  $P$ contains $\mathrm{id}=gg^{-1}$ and $\mathcal G$: for
$g\in\mathcal G$ pick $h\in\mathcal G$ with $h\ind{}g$; then $gh\in\mathcal G$ and
$g=(gh)h^{-1}\in P$.  $P$ is closed under inverses by \ref{H:inv}, and under products:
$(g_1g_2)(g_3g_4)=(g_5g_6)g_4=g_5(g_6g_4)$ is a product of two elements of $\mathcal G$ by
two applications of \cref{lem:threetwo}.  Hence $P$ is a subgroup containing
$\mathcal G$, so $P=G$.  $\Gamma$-invariance: $\sigma(g_1g_2)=\sigma g_1\,\sigma g_2$ and
$\mathcal G$ is $\Gamma$-invariant.
\end{proof}

\begin{proof}[Proof of \ref{hw:stab}]
Let $\gamma\in G$ and $h\in\mathcal G$ with $h\ind{}\gamma$.  By \ref{hw:two} write
$\gamma=g_1g_2$ with $g_i\in\mathcal G$.  The pair $(g_1,g_2)$ may be replaced by any
$\Gamma_{\acl(\Cb\gamma)}$-conjugate pair (its product is still $\gamma$, as such
automorphisms fix $\gamma$ by \cref{lem:germ-aut}), so by \cref{lem:fresh-conj} applied
to the code pair over $\acl(\Cb\gamma)$ we may assume
$g_1g_2\ind{\gamma}h$; together with $h\ind{}\gamma$ this gives $h\ind{}g_1g_2\gamma$
(\ref{R:trans}), in particular $h\ind{}g_1g_2$.  Then $g_2h\in\mathcal G$ by \ref{H:comp}, and
$\rk(g_2h/g_1)\ge\rk(g_2h/g_1g_2)=\rk(h/g_1g_2)=n$ (interalgebraicity of $g_2h$ and $h$
over $g_2$), so $g_2h\ind{}g_1$ and $\gamma h=g_1(g_2h)\in\mathcal G$ by \ref{H:comp}.
Moreover $\rk(\gamma h/\gamma)\ge\rk(\gamma h/g_1g_2)=\rk(h/g_1g_2)=n$ since $\gamma h$ and
$h$ are interalgebraic over $g_1g_2$ (which contains the code of $\gamma$); hence
$\gamma h\ind{}\gamma$.  For $h\gamma$ apply this to $\gamma^{-1}$ and $h^{-1}$
(\ref{H:inv}, $\Cb(h^{-1})=\Cb(h)$, $\Cb(\gamma^{-1})=\Cb(\gamma)$): $\gamma^{-1}h^{-1}\in\mathcal G$
and $\gamma^{-1}h^{-1}\ind{}\gamma^{-1}$, i.e.\ $h\gamma\in\mathcal G$ and $h\gamma\ind{}\gamma$.
\end{proof}

\begin{proof}[Proof of \ref{hw:rank}]
Let $\gamma\in G$ and pick $h\in\mathcal G$ with $h\ind{}\gamma$.  By \ref{hw:stab}
$h\gamma\in\mathcal G$, so $\rk(h\gamma)=n$.  Since $h\gamma$ and $\gamma$ are interalgebraic
over $h$, $\rk(\gamma)=\rk(\gamma/h)=\rk(h\gamma/h)\le\rk(h\gamma)=n$.  If $\rk(\gamma)=n$ then
$\rk(h\gamma/h)=n=\rk(h\gamma)$, i.e.\ $h\gamma\ind{}h$, so $h^{-1}\ind{}h\gamma$ and
$\gamma=h^{-1}(h\gamma)\in\mathcal G$ by \ref{H:inv}, \ref{H:comp}.  Conversely elements
of $\mathcal G$ have rank $n$.
\end{proof}

\begin{proof}[Proof of \ref{hw:action}]
The first claim is \cref{lem:eval}(2) with $F=\acl(\Cb\gamma)$.  The second is the
defining property of the composite $\gamma\circ\delta$ in \cref{lem:comp}, since
$y\ind{}\gamma\delta$ says that $y$ is generic over $\acl(\Cb\gamma\,\Cb\delta)$.
\end{proof}

\begin{lemma}[the key rank computation]\label{lem:keyrank}
Let $\gamma\in G$, $h\in\mathcal G$ and $w$ a tuple with $h\ind{}\gamma w$.  Then
$h\gamma\in\mathcal G$ and $h\gamma\ind{}w$; in particular, if $y\in V_{\Gen}$ and
$h\ind{}\gamma y$ then $y\in\operatorname{Dom}(h\gamma)$.
\end{lemma}
\begin{proof}
$h\gamma\in\mathcal G$ by \ref{hw:stab}.  As $h\gamma$ and $h$ are interalgebraic over
$\gamma$, $\rk(h\gamma/w)\ge\rk(h\gamma/w\gamma)=\rk(h/w\gamma)=n$.
\end{proof}

\begin{proof}[Proof of \ref{hw:space}]
\emph{$E$ is well defined.}  Let $\gamma,\gamma'\in G$, $y,y'\in V_{\Gen}$ and
$w:=\gamma\gamma'yy'$.  By \cref{lem:genset-basic}(1) there is at least one $h\in\mathcal G$
with $h\ind{}w$ (so the ``for some'' and ``for every'' forms are equivalent once the
condition is shown to be independent of $h$; a formalization should define $E$ by the
universal form).  Let $h,\tilde h\in\mathcal G$ both be independent from $w$.  By
\cref{lem:keyrank}, $y\in\operatorname{Dom}(h\gamma)$ and $y'\in\operatorname{Dom}(h\gamma')$, so
both sides of the defining equation are defined; likewise for $\tilde h$.  By
\cref{lem:genset-basic}(1) there is $\sigma\in\Gamma_{\acl(w)}$ with $\sigma h=\tilde h$;
$\sigma$ fixes the codes of $\gamma,\gamma'$ and the points $y,y'$, hence (\cref{lem:germ-aut},
\cref{lem:eval}(4)) transforms the equation for $h$ into the equation for $\tilde h$.
So the condition does not depend on $h$.

\emph{$E$ is an equivalence relation.}  Reflexivity and symmetry are clear.  For
transitivity, given $(\gamma,y)E(\gamma',y')$ and $(\gamma',y')E(\gamma'',y'')$ choose one
$h\in\mathcal G$ independent from all six objects and compare.

\emph{The action is well defined.}  Let $(\gamma,y)E(\gamma',y')$ and $\delta\in G$.
Choose $\tilde h\in\mathcal G$ with $\tilde h\ind{}\gamma\gamma'yy'\delta$.  By
\cref{lem:keyrank} (with $w=\gamma\gamma'yy'$), $\tilde h\delta\in\mathcal G$ and
$\tilde h\delta\ind{}\gamma\gamma'yy'$, so $\tilde h\delta$ is an admissible witness for
$(\gamma,y)E(\gamma',y')$: $(\tilde h\delta\gamma)(y)=(\tilde h\delta\gamma')(y')$.  Since
$\tilde h\ind{}\delta\gamma\,\delta\gamma'\,yy'$ as well (its rank over these is still $n$
because they lie in $\acl$ of the previous tuple), $\tilde h$ witnesses
$(\delta\gamma,y)E(\delta\gamma',y')$.  Associativity and unit law are inherited from $G$.

\emph{$\iota$ is injective.}  If $(\mathrm{id},y)E(\mathrm{id},y')$ then $h(y)=h(y')$ for
$h\in\mathcal G$ independent from $yy'$, and $y=h^{-1}(h(y))=h^{-1}(h(y'))=y'$ by
\cref{lem:comp} ($y,y'$ are generic over $\acl(\Cb h)=\acl(\Cb h^{-1})$).

\emph{Compatibility.}  If $y\ind{}\gamma$, choose $h\ind{}\gamma y\gamma(y)$; then
$y\ind{}\gamma h$ (\ref{R:trans}: $y\ind{}\gamma$ and $y\ind{\gamma}h$, the latter from
$h\ind{}\gamma y$), so $(h\gamma)(y)=h(\gamma(y))$ by \ref{hw:action}, i.e.\
$(\gamma,y)E(\mathrm{id},\gamma(y))$, which is $\gamma\cdot\iota(y)=\iota(\gamma(y))$.

\emph{Transitivity.}  Since $[\gamma,y]=\gamma\cdot\iota(y)$ it suffices to carry $\iota(y)$
to $\iota(y')$ for $y,y'\in V_{\Gen}$.  Choose $c\in V_{\Gen}$ with $c\ind{}yy'$.
The pairs $(y,c)$ and $(y',c)$ are independent generic pairs of $V$.  By \ref{H:trans}
and \cref{lem:genset-basic}(2), for any $g_0\in\mathcal G$ with $g_0\ind{}y$ the pair
$(y,g_0(y))$ is an independent generic pair; by \cref{lem:indep-conj} there is
$\sigma\in\Gamma$ with $\sigma(y,g_0(y))=(y,c)$, and then $g:=\sigma g_0\in\mathcal G$
satisfies $g\ind{}y$ and $g(y)=c$ (\cref{lem:eval}(4)).  Likewise there is $g'\in\mathcal G$
with $g'\ind{}c$ and $g'(c)=y'$.  By compatibility, $g\cdot\iota(y)=\iota(c)$ and
$g'\cdot\iota(c)=\iota(y')$, so $(g'g)\cdot\iota(y)=\iota(y')$.

\emph{Faithfulness.}  Suppose $\delta\in G$ fixes every element of $S$.  Choose
$y\in V_{\Gen}$ with $y\ind{}\delta$ and $h\in\mathcal G$ with $h\ind{}\delta y$.  From
$(\delta,y)E(\mathrm{id},y)$ we get $(h\delta)(y)=h(y)$.  Now $h\delta,h\in\mathcal G$ and
$y\ind{}h\delta$ (as $y\ind{}\delta$ and $h\ind{}\delta y$ give $y\ind{}\delta h$), so by
\cref{lem:eval}(5) $h\delta=h$ and $\delta=\mathrm{id}$.
\end{proof}

\begin{remark}[what ``connected'' and ``homogeneous space of rank $m$'' mean here]
\label{rem:connected}
Statement \ref{hw:rank} is the germ-level content of connectedness: the elements of
maximal rank form one $\Gamma$-orbit, and \ref{hw:stab} says that this orbit is
stabilised by left and right multiplication by independent elements.  The set $S$
contains (a copy of) $V_{\Gen}$, on which $G$ acts generically by \ref{hw:action}:
elements of $G$ move generic points of $V$ independent from them to generic points of
$V$ independent from them.  We define no rank on $S$ and make no claim about ``generic
elements of $S$''.  Every element of $S$ is $\gamma\cdot\iota(y)$ for some $\gamma\in G$ and
$y\in V_{\Gen}$ (by definition of $S$), but $y$ can \emph{not} in general be taken
independent from $\gamma$: for the group of translations of the affine line
($V=\mathbb A^1$, $\mathcal G$ the germs $y\mapsto y+a$ with $a$ transcendental, $G$ all
translations, $(\gamma,y)\,E\,(\gamma',y')$ iff $\gamma+y=\gamma'+y'$, so $S\cong K$ and
$\iota(V_{\Gen})$ is the set of transcendental elements) the class $[-t,t]=0$ has no
presentation $\gamma\cdot\iota(y)$ with $y\ind{}\gamma$, since by compatibility such a
presentation would equal $\iota(\gamma(y))\in\iota(V_{\Gen})$.  No later statement needs
such presentations.  In \cref{sec:constructibility} both $G$ and $S$ are shown to be
constructible quotients.
\end{remark}

\begin{remark}[comparison with the literature]
\Cref{thm:hw} is Theorem 5.1 of Bays' Hrushovski--Weil notes (= Lemma 5.2 of Bays'
geometric stability theory notes, which follow Pillay's \emph{Geometric Stability
Theory}, Chapter 5, \S1; Hrushovski's thesis, Theorem 3.1.1 together with the
Proposition on p.~37) written for germs of quasi-rational maps.  The only
``stability'' used is the uniqueness of generic points up to conjugacy
(\cref{lem:generic-char}, \cref{lem:indep-conj}) and the rank calculus.  The hypothesis
\ref{H:inv} replaces the chain-condition argument of the thesis, which is unnecessary in
the group-configuration application since inverses are available there by construction.
\end{remark}
